\section{Evaluation: The Appraisal of a Shared Life and Its Cost}
\label{sec:evaluation}

\noindent
The cycle turns to evaluation in its explicit form. The basal appraisal of \xref{sec:appraisal} runs continuously through every moment; the present moment is the deliberate, instrumented assessment undertaken at particular junctures, in which some aspect of the relation is brought to a structured and assessable form. This is the moment at which the operational concession (\xref{sec:bg-concession}) is most consequential, and the moment at which the positivist framework is most directly drawn upon, though, as will appear, every framework has its own evaluative practice and the frameworks do not agree on what evaluation is. The section first sets out the several evaluative practices, then develops, as its principal contribution, a quantitative assessment of the justice of a relation's value cycle.

\subsection{The several practices of evaluation}
\label{sec:ev-layers}

Evaluation is not one practice but several, drawn from the several frameworks, each assessing a different aspect in a different register.

The \emph{biological} appraisal is the continuous interoceptive and affective evaluation of \xref{sec:ap-bio}: the felt registration, beneath reflection, of how the relation stands. It is the most pervasive evaluation and the one conducted in the value nearest the real, and it is at the same time the least communicable and the least accountable, since it is unsymbolised. It is the ground the other evaluations structure, and not a competitor to them.

The \emph{decision-theoretic} appraisal is the belief updating of \xref{sec:ap-pomdp}: the revision, on the evidence of the lived fabric, of each party's distribution over the state of the relation. It is evaluation in the register of probability rather than verdict, and it carries the recognition that the state assessed is hidden and the assessment therefore irreducibly uncertain.

The \emph{psychological} appraisal is the structured self-report of appraisal theory and its instruments: the assessment of the relation through the parties' reports of their own appraisals, as in the measurement of relationship satisfaction. It renders the basal appraisal partially communicable by structuring it into reportable form.

The \emph{positivist} evaluation is the measurement proper: the assessment of the relation by validated instruments yielding comparable and accumulable quantities. Its instruments are many, and they measure distinct things. Life-satisfaction and affect scales measure subjective well-being; quality-of-life instruments measure the conditions of a life; relationship-satisfaction and adjustment scales measure the parties' assessment of the relation; eudaimonic and self-determination instruments measure reported growth, autonomy, and the sense of an unfolding life. Representative instruments are, respectively, the satisfaction-with-life and affect scales \citep{diener1985,watson1988}, quality-of-life assessments \citep{whoqol1998}, dyadic-adjustment scales \citep{spanier1976}, psychological well-being scales \citep{ryff1989}, and measures grounded in self-determination theory \citep{ryan2000,deci1985sdt}; the eudaimonic turn in this measurement is represented by \citet{seligman2011}. Each is a structured proxy, validated for its reliability and its correlation with other measures, and each is, within its operational limits, a real gain: public, comparable across relations and across time, and accumulable into a body of findings. The positivist evaluation is the developed form of explicit assessment, and the paper neither disparages it nor mistakes it for the appraisal of the real value it proxies.

The \emph{phenomenological} evaluation is the thick description and the hermeneutic assessment of the shared life: the close interpretive characterisation of how a relation stands, conducted in narrative rather than number, and answerable to the lived meaning of the relation rather than to a comparable metric. It is evaluation that resists the structuring the positivist evaluation requires, on the ground that the structuring discards what is to be assessed, and it stands to the positivist evaluation as the interpretation of a text stands to the counting of its words.

These five practices do not agree on what evaluation is, and the disagreement is not to be resolved by the victory of one. They are evaluative phases, in the sense of \xref{sec:bg-phases}: the positivist measures and is silent on what it cannot structure; the phenomenologist interprets and resists the metric; the decision-theorist quantifies uncertainty and the psychologist structures self-report; the biological appraisal grounds them all in a value none of them reaches. A relation is evaluated well when the practices are drawn upon together and none is mistaken for the whole, and most particularly when the structured measures are read as proxies and not as the appraisal of the real.

\input{sections/ev_survey.tex}

\subsection{The quantitative assessment of justice}
\label{sec:ev-justice}

The paper now develops, as the principal contribution of this moment, a quantitative assessment directed at the condition of justice in the normative specification of \xref{sec:cog-spec}. The condition of justice is that the value generated in the shared life circulates and returns to those who generate it, no party being reduced to a source from which value is extracted without return. The assessment of this condition admits of a structured proxy that makes visible a feature the satisfaction measures cannot, the systematic maldistribution of return that constitutes exploitation within an apparently satisfying relation. The assessment is offered under the operational concession throughout: it is a structured proxy for the justice of a relation's value cycle, operationally valuable and not a measurement of the relation's real worth.

\subsubsection{Expected recyclability}
\label{sec:ev-recyc}

The first quantity is the expected recyclability of value in the relation: the probability that a value generated by one party is met, within the relation, by a response that regenerates value in turn, rather than being absorbed without return. A contribution to a shared life, an act of care, an effort of attention, a piece of the relation's sustaining labour, may be met in one of three ways. It may be met by a response that itself generates value, an acknowledgement, a reciprocation, a deepening, so that the contribution sets a further turn of the cycle in motion; this is the recyclable case. It may be met by a response that neither regenerates nor depletes, a neutral absorption in which the contribution is received and nothing further follows; this is the lossy case. Or it may be met by a response that depletes, a contribution received and answered with a withdrawal or a demand for more, so that the act of giving leaves the giver with less than before; this is the extractive case. The expected recyclability is the probability, characteristic of the relation, that a contribution falls into the first case rather than the second or the third.

So defined, recyclability is the proxy proper to the condition of the cycle's life. A cycle sustains itself when the value it regenerates, turn upon turn, at least balances the value it loses to neutral absorption and extractive response; and whether it does so is governed by the recyclability. There is a threshold: above a critical recyclability, contributions regenerate faster than they are lost, and the cycle of value is self-sustaining or growing; below it, contributions are lost faster than they regenerate, and the cycle winds down, its sustaining labour meeting diminishing response until the parties cease to contribute and the relation decays toward the extinction of \xref{sec:regimes}. The condition is recognisable in its phenomenology. A relation of high recyclability is one in which giving feels generative, in which an effort made is met by a response that makes further effort worth making; a relation of low recyclability is one in which giving feels thankless, in which contributions disappear into a partner who receives without regenerating, and in which the giver, over time, contributes less because the contributions return nothing. The decline of recyclability below its threshold is the quantitative form of a relation's running down, and it is to be distinguished, as the next quantities require, from the maldistribution of return among the parties, which a relation of adequate aggregate recyclability may still exhibit.

\subsubsection{Return relative to a coupling-based baseline}
\label{sec:ev-return}

The second quantity concerns the distribution of return among the parties, assessed against a baseline set by the coupling of the relation. The intuition the quantity formalises is that the degree of return that is just is not fixed but depends on how deeply the parties are coupled. Coupling, here, is the degree to which the parties' goods are bound together, the extent to which what benefits one benefits the other and what the relation generates is held in common rather than apportioned. In a loosely coupled relation, say an association of convenience, the parties' goods are largely separate, and a contribution by one may justly yield most of its return outside the relation, to the contributor's separate good; little is owed back through the relation, because little was staked in common. In a deeply coupled relation, a shared life fully joined, the parties' goods are bound together, and a contribution by one is made into a common store from which return is owed back to its generator through the relation; much is owed back, because much was staked in common. The coupling thus sets a baseline expected return: the share of the value a party generates that, given the depth of the coupling, ought to circulate within the relation and return to that party rather than be absorbed elsewhere.

The baseline is what makes the assessment of return a normative and not merely a descriptive matter. To observe that a party's contributions return little to that party is not yet to find an injustice, for in a loosely coupled relation little return is owed. The injustice is the shortfall of actual return below the baseline the coupling sets: a party in a deeply coupled relation, who has staked a shared life in common and whose contributions therefore ought to return to that party through the relation, but whose actual return falls systematically below that baseline, is a party from whom value is extracted under the very depth of coupling that ought to have returned it. This is the characteristic structure of exploitation within intimacy, and it is invisible to an assessment that does not set the baseline by the coupling: the deeply coupled party who is under-returned has given most and is owed most, and it is precisely the depth of the coupling, read without the baseline, that can be made to look like generosity freely given rather than value unjustly retained. The assessment compares the actual return to each party against the coupling-based baseline, and reads a systematic shortfall as extraction.

\subsubsection{Skewness as the statistical fingerprint of injustice}
\label{sec:ev-skew}

The third quantity is the principal instrument of the assessment. Consider the distribution, across the parties and across time, of value returned relative to value generated. In a just relation, this distribution is, relative to the coupling-based baseline, approximately symmetric: returns are distributed among the generators in proportion to their generation, and no generator is systematically consigned to the low-return tail. In an exploitative relation, the distribution is skewed: value returns disproportionately to one party, while another party, characteristically the party who bears the larger share of the relation's sustaining labour, generates value that is systematically under-returned and is consigned to the tail of the distribution.

\begin{claim}[Skewness of the return distribution as the statistical fingerprint of injustice]
The injustice of a relation's value cycle has a statistical signature in the skewness of the distribution of returned value relative to the coupling-based baseline. A just cycle returns value to its generators in proportion to their generation and to the coupling, yielding a distribution that is approximately symmetric about the baseline. An exploitative cycle returns value disproportionately to one party and consigns another systematically to the low-return tail, yielding a distribution skewed away from the under-returned party. The skewness of the return distribution is therefore a structured proxy for the injustice of the cycle, and it makes visible a maldistribution that aggregate measures of satisfaction conceal: a relation may exhibit high mean satisfaction while its return distribution is severely skewed, the high mean being the satisfaction of the over-returned party and the skew the signature of the other's exploitation. The skewness is read, under the operational concession, as a proxy for injustice and not as its measurement; a symmetric distribution is necessary but not sufficient for justice, and a skewed one is evidence of, and not identical to, exploitation.
\end{claim}

\noindent
The skewness measure has the methodological merit of being orthogonal to the satisfaction measures, and so of detecting what they systematically miss. An aggregate or mean measure of relational satisfaction is raised by the satisfaction of the over-returned party and is insensitive to the maldistribution beneath it; a relation can score well on mean satisfaction precisely because the party it favours is well satisfied, while the party it exploits is consigned to the tail. The skewness reads the shape of the distribution rather than its mean, and so detects the maldistribution the mean conceals. This is the formal counterpart, in the assessment of a relation, of the prior paper's thesis that a satisfaction purchased by the foreclosure of another is no eudaimonia: the mean satisfaction is high, and the relation is nonetheless unjust, and the injustice is legible in the skew.

The mechanism by which one party comes to occupy the low-return tail deserves to be made concrete, because it is the recurring structure the measure is designed to detect. A relation distributes, among other things, its sustaining labour: the work of attention, of care, of maintaining the household and the relationship itself, the emotional and reproductive labour without which the shared life does not continue. The systematic under-recognition of this labour as labour, and its uncompensated extraction, is documented by \citet{hochschild1983} and analysed in the economics of care by \citet{folbre2001}. This labour generates value, much of it, since the shared life it sustains is the medium of the relation's whole flourishing. When this labour is borne disproportionately by one party, and when the value it generates returns disproportionately to the other, the bearer of the sustaining labour is the party consigned to the tail: she generates a large share of the relation's value and receives a small share of its return, and the gap is the measure of her exploitation. The structure is the more durable for being concealed by its own ideology, in which the sustaining labour is figured as freely given, as natural to the one who bears it, or as no labour at all, so that its under-return is read not as extraction but as the order of things. The skewness measure is constructed to make this structure legible against that concealment: it does not ask whether the bearer of the labour reports satisfaction, which she may, having been habituated to expect no return, but whether the value she generates returns to her, which, in the exploitative case, it does not.

A contrast makes the orthogonality to the mean vivid. Consider two relations that score identically on a measure of mean satisfaction. In the first, both parties contribute to the sustaining labour, the value each generates returns to each in proportion to the coupling, and the return distribution is approximately symmetric; the shared mean satisfaction is the mark of a flourishing that both enjoy. In the second, one party bears almost all the sustaining labour and receives little return, while the other receives most of the relation's value and bears little of its cost; the over-returned party is highly satisfied, the under-returned party reports a satisfaction depressed but not absent, having adjusted her expectations to her position, and the mean of the two is, by construction, equal to the mean of the first relation. The mean satisfaction does not distinguish these relations; the skewness does. The first has a symmetric return distribution and the second a severely skewed one, and the skew of the second is the signature of an exploitation that the equal means conceal. The measure earns its place exactly here: it sees, in two relations a satisfaction survey would rank together, the difference between a flourishing both parties share and a flourishing one party extracts from the other.

\subsubsection{A model for the test, and a hypothesis}
\label{sec:ev-model}

The skewness measure presupposes a model, since a distribution of return cannot be assessed for its skewness without a model of how return is generated, against which the observed distribution is compared. The paper indicates the form of the required model without claiming to have validated it, in keeping with the concession. The relation is represented as a network whose nodes are the generators, the two parties and the coupled others, characteristically family, whose value the relation also circulates, and whose edges carry the couplings, the weighted channels along which value propagates from one generator to another. Value generated at a node does not remain there; it propagates along the edges, a portion returning to its origin, a portion passing to coupled others, a portion lost from the network altogether. The propagation is modelled as a branching process on this network, in the manner of the prior series' formalisation of generative recursion: the mathematics of such propagation, in which quantities at nodes generate quantities at successors and the probability of eventual return is the object of interest, is the classical theory of branching processes \citep{harris1963}; value at a node produces, in expectation, value at adjacent nodes in proportion to the edge weights, each such product producing further value in turn, and the expected return to a generator is the probability, accumulated over all the paths by which value may travel and return, that value originating at that generator's node finds its way back to it. The coupling of the relation enters twice, as the topology of the network, which generators are joined to which, and as the weights on the edges, how strongly; and together these fix, for each generator, the coupling-based baseline of \xref{sec:ev-return}, the expected return that the coupling structure would yield were propagation unbiased, neither favouring nor disfavouring any generator beyond what the couplings themselves dictate.

On such a model a hypothesis test becomes statable, and its terms can be given exactly. Let the normalised return to a generator be its actual expected return divided by its coupling-based baseline, so that a value of one is the return the coupling alone would yield, a value above one a generator favoured beyond its coupling, and a value below one a generator disfavoured below it. The null hypothesis is that return is unbiased: that the normalised return is, in expectation, equal across generators, so that the distribution of normalised return is symmetric about its centre and no generator is systematically favoured or disfavoured beyond the coupling. The alternative is that return is systematically biased: that some generator, characteristically the bearer of the sustaining labour, has a normalised return whose expectation lies significantly below the others, so that the distribution is significantly skewed against that generator. The test statistic is the skewness of the observed distribution of normalised return, and the test asks whether an observed skew lies within the range compatible with the null of unbiased propagation or is large enough to reject it in favour of systematic extraction. The paper advances the form of this test, and not a validated instrument. The model requires assumptions that have not here been discharged, on the propagation probabilities, on the independence of the branching, and on the stationarity of the couplings over the period of estimation; the estimation of returns in an actual relation, the assignment of value to acts of care and its tracing through the network, is a substantial empirical problem the paper does not solve; and the result, were it obtained, would be a structured proxy for injustice under the concession, evidence of extraction and not its measurement. The contribution is the form: that the justice of a relation's value cycle is, in principle and under the concession, a testable hypothesis about the skewness of a return distribution against a coupling-based baseline, and not a matter on which only impression can be brought to bear.

\subsection{The operationalisation of return, and a worked illustration}
\label{p11:sec:ev-operational}

The assessment of justice requires that the abstract notion of the return of value be operationalised into a quantity that could, in principle, be estimated; this subsection sets out the operationalisation and illustrates it on a deliberately simple example, under the operational concession throughout.

The operationalisation proceeds in four steps. First, the contributions of each party over a period are identified: the discrete acts of generation, the pieces of sustaining labour, the cared-for occasions, that each party contributes to the shared life. Second, to each contribution is assigned a value, in a common unit, proxying the value it generates for the shared life; the assignment is the most contestable step, and the concession applies most stringently here, since the value of an act of care is exactly the unsymbolised quantity the proxy cannot reach, and the assignment is an operational stipulation and not a measurement. Third, for each party the return is identified: of the value generated in the shared life, the portion that returns to that party, whether directly or through the relation's circulation. Fourth, the return to each party is normalised by the coupling-based baseline of \xref{p11:sec:ev-return}, the return the coupling structure would yield under unbiased propagation, to give the normalised return on which the assessment of skewness is conducted.

A worked illustration, with stipulated numbers, shows the quantities in operation; the numbers are illustrative and carry no empirical claim. Consider a deeply coupled relation of two parties, $A$ and $B$, over a period in which the shared life generates value, in the common unit, of 100. Suppose $B$ bears the larger share of the sustaining labour and generates 60 of this value, while $A$ generates 40. In a just cycle, given the depth of the coupling, the return to each would be approximately proportional to generation: $B$ would receive a return near 60 and $A$ a return near 40, the normalised returns of both being near 1, and the distribution of normalised return approximately symmetric. Suppose instead that the relation returns value disproportionately to $A$: of the 100 generated, $A$ receives a return of 65 and $B$ a return of 35. The normalised return to $A$ is then $65/40 \approx 1.6$ and to $B$ is $35/60 \approx 0.58$. The distribution of normalised return is now skewed: $A$ sits well above the baseline of 1, and $B$ well below it, in the low-return tail. The skewness of the normalised-return distribution, negative for $B$, is the statistical fingerprint of \xref{p11:sec:ev-skew}, and it is legible here despite the relation's possibly high mean satisfaction, since $A$, the over-returned party, may report great contentment, and $B$, habituated to bearing the sustaining labour for little return, may report a satisfaction only moderately depressed, so that the mean of the two is unremarkable while the skew is severe. The illustration shows the quantities the assessment computes; it does not show that they can be estimated easily in an actual relation, which, as \xref{p11:sec:ev-model} stated, remains a substantial empirical problem the paper does not solve.


\subsection{The evaluation of the unfolding of each}
\label{sec:ev-dynamics}

The condition of justice admits the structured proxy just developed. The eudaimonic condition, the unfolding of each party's own dynamics, is the most resistant to structuring, because it lies nearest the unsymbolised value the proxy cannot reach, and the paper records its resistance rather than overcoming it. Two proxies are nonetheless available and worth stating, under the concession.

The first is the rate at which each party's intended undertakings are in fact accomplished: the proportion of what a party sets out to be and to do that the party, within the shared life, achieves. The intended undertakings of a party are the projects, the developments, the becomings that party has reason to value and sets out upon, the unfolding of that party's own dynamics in concrete form: a body of work, a craft, a friendship, a course of study, a way of living the party means to grow into. The proxy tracks, over time, what proportion of these a party accomplishes rather than abandons. A shared life that sustains the unfolding of each is one in which each party's intended undertakings are, over time, accomplished at a rate that does not systematically decline, the relation making room for and conducing to the development each sets out upon; a shared life that forecloses one party's unfolding is one in which that party's intended undertakings are systematically unaccomplished, the development that party set out upon stalling and then ceasing to be attempted. The rate of accomplishment is a coarse proxy, since not every intended undertaking ought to be accomplished and not every abandonment is a foreclosure, but a systematic decline in one party's rate, sustained over the seasons of a shared life, is a structured signature of a unfolding being foreclosed.

The second proxy refines the first by attributing the shortfall, and it is the more telling of the two. Of a party's intended undertakings that are not accomplished, the assessment distinguishes those not accomplished for reasons extraneous to the relation, the ordinary attrition of projects that any life suffers, from those not accomplished because the shared life foreclosed them: the undertakings a party set aside, deferred, or abandoned on account of the relation, because the relation's demands left no room for them, because the sustaining labour the party bore consumed the time and the energy they required, or because the party's development was tacitly understood, by both parties, to be the one that would yield. It is this attributed shortfall, and not the gross rate of abandonment, that proxies foreclosure. A party may accomplish few of their intended undertakings for reasons that have nothing to do with the relation, and a relation is not thereby indicted; but a party who systematically abandons their undertakings on account of the relation, while the other party's undertakings proceed, is a party whose unfolding the relation is foreclosing in favour of the other's. A systematically high proportion of intended undertakings foreclosed on account of the relation, concentrated in one party, is the proxy for the foreclosure of that party's unfolding, and it couples directly to the assessment of justice. The party whose unfolding the relation forecloses is characteristically the party consigned to the low-return tail of \xref{sec:ev-skew}, for the same sustaining labour that returns little value to its bearer is what consumes the room her own undertakings would have required; the two proxies, the skew of return and the concentration of foreclosed undertakings, are two structured signatures of one exploitation, read in the register of value returned and in the register of development foreclosed.

\begin{claim}[The foreclosure of unfolding as a proxy coupled to the skew of return]
The foreclosure of a party's unfolding admits a structured proxy in the proportion of that party's intended undertakings that are abandoned on account of the relation, and this proxy is coupled to the skew of the return distribution. The party whose unfolding the relation forecloses is characteristically the party from whom value is extracted without return, so that a concentration of foreclosed undertakings in one party and a skew of the return distribution against the same party are two signatures of a single injustice, the reduction of that party to the sustaining fuel of a relation that returns to it neither value nor the freedom to unfold. The proxies are read under the concession, as structured signatures of an exploitation they evidence and do not measure, and they leave the unsymbolised remainder of each party's unfolding beyond their reach.
\end{claim}

\subsection{The poetic condition, and the limit of evaluation}
\label{sec:ev-poetic}

The third condition of the normative specification, that the cycle be sustained as a living and open movement rather than settled into a closed repetition, is the least amenable to structuring, and the paper records only the direction a proxy would take. The condition is that the relation's recurrences, its habits, its shared signs, its rituals, continue to accumulate an irreducible remainder on each traversal, the positive holonomy of the prior paper, rather than emptying into a repetition that returns nothing new. A proxy would assess whether the relation's recurrences are generative, whether a shared practice traversed again yields something not contained in its prior traversals, or have become a closed repetition traversed without gain. The paper does not develop a measure of this, and holds that it is at best semi-quantitative, since the accumulation of an irreducible remainder is close to the unsymbolised value that no structuring reaches. The poetic condition marks the limit at which the evaluation of a shared life passes beyond the structured proxy and returns to the basal and unsymbolised appraisal from which it was abstracted.

\subsection{The standing of the whole evaluation}
\label{sec:ev-standing}

The evaluation developed in this section, in all its practices and proxies, stands under the operational concession, and the concession is to be restated as the section closes, because the next moment of the cycle, reflection, turns precisely on it. Every structured measure of this section, the satisfaction scales, the recyclability, the skew of return, the rate and attribution of accomplished undertakings, is a structured proxy for a value it does not reach, operationally valuable and incomplete. The measures are public where the basal appraisal is private, comparable where it is singular, accumulable where it is evanescent, and these are real gains that warrant the measures' use. The measures are, at the same time, structured symbolisations of an appraisal whose object is unsymbolised, and they discard exactly the remainder in which the relation's real value consists. To use them is to accept this exchange knowingly, for the operational gain it brings; to mistake them for the appraisal of the real value is the error the next moment of the cycle, reflection, exists to correct.
